\documentclass[a4paper, 10pt, ]{article}

\input{../../KUT000/preamble_v1.0.tex}

% -----------------------------------------------------------------------------

\def\oznacenieCelku{Kolekcia učebných textov}

% -----------------------------------------------------------------------------


\def\KUTporadoveCislo{021}

% \def\oznacenieVerzie{v1.0}
\def\oznacenieVerzie{\phantom{v1.0}}

\def\mesiacRok{december 2025}

\def\authorslabel{MT, RJ}






% -----------------------------------------------------------------------------

\begin{document} 

% -----------------------------------------------------------------------------
% Uvodny nadpis

\noindent
\parbox[t][18mm][c]{0.3\textwidth}{%
\raisebox{-0.9\height}{%
\phantom{.}
\includegraphics[height=18mm]{../../KUT000/URKFEIlogo_v0.1.pdf}%
}%
}%
\parbox[t][18mm][c]{0.7\textwidth}{%
\raggedleft

\sffamily
\fontsize{16pt}{18pt}
\fontseries{sbc}
\selectfont

\noindent
\textcolor[rgb]{0.75, 0.75, 0.75}{\textls[25]{\oznacenieCelku}}
}%

\noindent
\parbox[t][16mm][b]{0.5\textwidth}{%
\raggedright

\color{Gray}
\sffamily

\fontsize{12pt}{12pt}
\selectfont
\mesiacRok

\fontsize{6pt}{10pt}
\selectfont
github.com/OkoliePracovnehoBodu/KUT

\fontsize{8pt}{10pt}
\selectfont
\authorslabel




}%
\parbox[t][16mm][b]{0.5\textwidth}{%
\raggedleft

\sffamily

\fontsize{6pt}{6pt}
\selectfont

\textcolor[rgb]{0.68, 0.68, 0.68}{\oznacenieVerzie}


\fontsize{14pt}{14pt}
\selectfont

\bfseries


\includegraphics[height=12pt]{../../KUT000/KUT_logo_v0.1.pdf}%
{%
\textls[-50]{\KUTporadoveCislo}
}%
}%

% -----------------------------------------------------------------------------



\vspace{6mm}

% ---------------------------------------------
\sffamily
\bfseries
\fontsize{18pt}{21pt}
\selectfont

\begin{flushleft}
    AeroShield:\\ návrh gray-box modelu bez vstupu
\end{flushleft}

\bigskip

% -----------------------------------------------------------------------------
\normalsize
\normalfont
% -----------------------------------------------------------------------------

\lstset{style=mystyle}










\noindent
\lettrine[lines=1, nindent=1pt, loversize=0.0]{C}{ieľom} 
textu je opis a matematického modelu zodpovedajúceho pohybu ramena laboratórneho zariadenia AeroShield. Ide o~model založený na fyzikálnych zákonoch a~z~hľadiska teórie systémov ide o model sivej skrinky (angl. \textit{grey-box model}).

{\color{Red}
Nasledujúce informácie nadväzujú na opis laboratórneho zariadenia AeroShield uvedený v~dokumente:
\medskip

\noindent
\begin{tabular*}{\textwidth}{ @{} >{\sffamily}p{2.0cm} @{\extracolsep{\fill}} p{11cm}<{\raggedright}}

    \href{run:../../KUT020/TeX/KUT020.pdf}{KUT020} & \href{run:../../KUT020/TeX/KUT020.pdf}{Lab. zariadenie AeroShield: orientačný prehľad} \\ \addlinespace[3pt]  

\end{tabular*}
}

\medskip


{\color{Red}
\section{Východiskové predpoklady}

\begin{itemize}[leftmargin=0pt, labelsep=3mm, itemsep=0pt]
    \item Uvažuje sa zariadenie ako pasívne kyvadlo bez pohonu.
\end{itemize}

}


\section{Pohybová rovnica}

Uvažujme ideálne tuhé teleso, ktoré je zavesené a môže sa otáčať okolo pevnej osi, ktorá je kolmá na rovinu jeho pohybu. Takéto teleso sa nazýva fyzikálne kyvadlo (obr.~\ref{kyv_itt_f1}). 

\begin{center}

    \vspace{-1mm}

    \vbox{%
    \makebox[\textwidth][c]{%
    \input{../SVG/kyv_itt_f1.pdf_tex}
	}

    \vspace{-6mm}

	\figcaption{Všeobecné fyzikálne kyvadlo}
    \label{kyv_itt_f1}
    }
	

    \vspace{-1mm}

\end{center}

Poloha ťažiska telesa je určená uhlom $\varphi$ [rad] a dĺžkou $L_T$ [m]. Uhol $\varphi$ je uhol,  ktorý zviera zvislá priamka prechádzajúca osou otáčania s~priamkou spájajúcou os otáčania a~ťažisko telesa. Dĺžka $L_T$ je vzdialenosť ťažiska od osi otáčania. 

Uhol $\varphi$ sa mení v~čase $t$ [s], označme $\varphi(t)$. Nulová poloha kyvadla, teda uhol $\varphi(t) = 0$ , je keď je kyvadlo v~zvislej polohe smerom dole.

Hmotnosť telesa je $m_k$ [kg], gravitačné zrýchlenie je $g$ [m/s$^2$]. 



\subsection{Potenciálna energia}

Vychýlením kyvadla z~nulovej polohy sa mení jeho potenciálna energia $U$ [J]. Vo všeobecnosti potenciálna energia je 
\begin{equation}
    U =  m_k \  g \  \left( L_T \ (1 -\cos(\varphi)) \right) 
\end{equation}
Z hľadiska pohybových rovníc sa zaujímame o~zmenu potenciálnej energie, typicky sa uvažuje vyjadrenie potenciálnej energie v tvare
\begin{equation}
    U = - \  m_k \  g \   L_T \  \cos(\varphi) 
\end{equation}




\subsection{Kinetická energia}

Kinetická energia rotačného pohybu ideálne tuhého telesa vo všeobecnosti je 
\begin{equation}
    K = \frac{1}{2} \  I_k \  \omega^2
\end{equation}
kde $I_k$ [kg$\ $m$^2$] je moment zotrvačnosti telesa vzhľadom na os otáčania a $\omega$ [rad/s] je uhlová rýchlosť telesa. V tomto prípade uhlová rýchlosť $\omega$ je 
\begin{equation}
    \omega(t) = \frac{\text{d}}{\text{d}t} \varphi(t) =   \dot\varphi(t)
\end{equation}
a teda kinetická energia je
\begin{equation}
    K = \frac{1}{2} \  I_k \  \dot\varphi^2
\end{equation}


\subsection{Lagrangeova funkcia a Lagrangeova rovnica}

Lagrangeova funkcia $L$ [J] je definovaná ako rozdiel kinetickej a potenciálnej energie
\begin{equation}
    L = K - U
\end{equation}
V uvažovanom prípade teda
\begin{equation}
    L = \frac{1}{2} \  I_k \  \dot\varphi^2 + m_k \  g \   L_T \  \cos(\varphi)
\end{equation}

Lagrangeova rovnica je vo všeobecnosti v tvare
\begin{equation} \label{lagrangeovaRovnica}
    \frac{\text{d}}{\text{d}t} \left( \frac{\partial L}{\partial \dot{q}} \right) - \frac{\partial L}{\partial q} = 0
\end{equation}
kde $q$ uvažujeme ako stupeň voľnosti pohybu systému (je možné voľne zvoliť jeho začiatočné podmienky/hodnoty) a táto rovnica tak vo všeobecnosti predstavuje pohybovú rovnicu systému. 

V uvažovanom prípade má systém jeden stupeň voľnosti, uhol $\varphi(t)$. Lagrangeova rovnica v tomto prípade je
\begin{equation}
    \frac{\text{d}}{\text{d}t} \left( \frac{\partial L}{\partial \dot{\varphi}} \right) - \frac{\partial L}{\partial \varphi} = 0
\end{equation}

Vyjadrime
\begin{equation}
    \frac{\partial L}{\partial \dot{\varphi}} = \frac{\partial}{\partial \dot{\varphi}} \left( \frac{1}{2} \  I_k \  \dot\varphi^2 + m_k \  g \   L_T \  \cos(\varphi) \right) = I_k \  \dot\varphi
\end{equation}
ďalej
\begin{equation}
    \frac{\text{d}}{\text{d}t} \left( \frac{\partial L}{\partial \dot{\varphi}} \right) = \frac{\text{d}}{\text{d}t} \left( I_k \  \dot\varphi \right) = I_k \  \ddot\varphi
\end{equation}
a nakoniec
\begin{equation}
    \frac{\partial L}{\partial \varphi} = \frac{\partial}{\partial \varphi} \left( \frac{1}{2} \  I_k \  \dot\varphi^2 + m_k \  g \   L_T \  \cos(\varphi) \right) = - m_k \  g \   L_T \  \sin(\varphi)
\end{equation}
Konkrétna Lagrangeova rovnica teda je
\begin{equation} \label{pohybovaRovnicaKyvadla}
    I_k \  \ddot\varphi(t) + m_k \  g \   L_T \  \sin(\varphi(t)) = 0
\end{equation}
a z fyzikálneho hľadiska ide o pohybovú rovnicu uvažovaného fyzikálneho kyvadla.



\subsection{Trenie}

V uvedenej Lagrangeovej rovnici \eqref{lagrangeovaRovnica} nie sú vo všeobecnosti uvažovné tzv. nepotenciálové sily (pravá strana rovnice je nulová). Vzhľadom na uvažované fyzikálne kyvadlo ide typicky o externý moment sily (pohon) a moment sily spôsobený trením.

Principiálne sa tu bude uvažovať moment sily $\tau_f$ [N m] spôsobený trením daný vzťahom
\begin{equation}
    \tau_f = F_{c} \left| \dot \varphi \right|^{z} \text{sgn}(\dot \varphi) 
\end{equation}
kde $\dot \varphi$ [rad/s] je uhlová rýchlosť, $F_{c}$ [N m/(rad/s)$^z$] je koeficient trenia a $z$ je exponent určujúci typ trenia. Pre $z = 0$ ide o suché trenie, pre $z = 1$ ide o viskózne trenie.

Uvažujme suché a viskózne trenie, potom pohybová rovnica \eqref{pohybovaRovnicaKyvadla} má tvar
\begin{equation} \label{pohybovaRovnicaKyvadlaTrenie}
    I_k \  \ddot\varphi(t) + m_k \  g \   L_T \  \sin(\varphi(t)) = - \tau_0 - \tau_1
\end{equation}
kde $\tau_0$ je moment sily spôsobený suchým trením a $\tau_1$ je moment sily spôsobený viskóznym trením.


\subsubsection{Suché trenie}

Nech konkrétny model suchého trenia je nasledovný
% \begin{subequations}
% \begin{align}
%     &\text{ak}  \quad \dot \varphi \neq 0 :&     \tau_0 &= F_{c0} \ \tanh(hc\cdot \dot \varphi) \label{sucheNormalMode} \\
%     &\text{ak}  \quad \dot \varphi = 0 \quad \text{a} \quad \varphi(0) \neq 0 :&  \tau_0 &= F_{c0} \ \tanh(hc\cdot \dot \varphi) \label{sucheInitMode} \\
%     &\text{inak} :&  \tau_0 &= - m g L_T \sin(\varphi) \label{sucheStuckMode}
% \end{align}
% \end{subequations}
% \begin{subequations}
% \begin{align}
%     \tau_0 &= F_{c0} \ \tanh(hc\cdot \dot \varphi) & \text{ak}   \ \dot \varphi &\neq 0  \label{sucheNormalMode} \\
%     \tau_0 &= F_{c0} \ \tanh(hc\cdot \dot \varphi) & \text{ak}   \ \dot \varphi &= 0 \quad \text{a} \quad t = 0 \quad \text{a} \quad \varphi(0) \neq 0  \label{sucheInitMode} \\
%     \tau_0 &= - m_k g L_T \sin(\varphi) &  (\text{inak})\  \text{ak}   \ \dot \varphi &= 0    \label{sucheStuckMode}
% \end{align}
% \end{subequations}
% \begin{subequations}
\begin{align}
    \tau_0 &= 
    \begin{cases} 
        F_{c0} \ \tanh(hc\cdot \dot \varphi) & \text{ak}   \ \dot \varphi \neq 0  \\
        F_{c0} \ \tanh(hc\cdot      \varphi) & \text{ak}   \ \dot \varphi = 0 \quad \text{a} \quad \left|\tau_s\right| \geq F_{c0} \\
        \tau_s &  \text{ak}  \ \dot \varphi = 0 \quad \text{a} \quad \left|\tau_s\right| < F_{c0}
    \end{cases}         
\end{align}
% \end{subequations}
kde $F_{c0}$ [N m] je koeficient suchého trenia a $hc$ je konštanta určujúca strmosť prechodu z~jedného smeru momentu sily do druhého smeru v okolí nulovej rýchlosti pričom použitie funkcie $\tanh()$ zabezpečuje spojitosť na rozdiel od použitia funkcie $\text{sgn}()$. 
% Vzťahom \eqref{sucheNormalMode} je dané suché trenie pri normálnom pohybe kyvadla (nenulová uhlová rýchlosť), vzťahom \eqref{sucheInitMode} je ošetrený prípad, keď je kyvadlo v začiatočnej polohe mimo nulovej polohy pretože inak vzťah \eqref{sucheStuckMode}, ktorý určuje veľkosť momentu sily suchého trenia pri nulovej uhlovej rýchlosti, by spôsobil, že kyvadlo by sa vôbec nepohlo.
Moment sily $\tau_s$ [N m] predstavuje súčet všetkých momentov síl pôsobiacich keď platí $\dot \varphi = 0$ okrem momentu sily spôsobeného suchým trením. V tomto prípade
\begin{equation}
    \tau_s = - m_k g L_T \sin(\varphi)
\end{equation}


\subsubsection{Viskózne trenie}
Nech konkrétny model viskózneho trenia je nasledovný
\begin{align}
    \tau_1 = F_{c1} \ \dot \varphi
\end{align}
kde $F_{c1}$ [N m/(rad/s)] je koeficient viskózneho trenia.




\subsection{Pohybová rovnica: zhrnutie}

Výstupnou veličinou uvažovaného systému je $\varphi(t)$, dynamika je daná diferenciálnou rovnicou
\begin{subequations}
\begin{equation} 
    I_k \  \ddot\varphi(t) + m_k \  g \   L_T \  \sin(\varphi(t)) = - \tau_0 - \tau_1
\end{equation}
pričom
\begin{align}
    \tau_s &= - m_k g L_T \sin(\varphi) \\
    \tau_0 &= 
    \begin{cases} 
        F_{c0} \ \tanh(hc\cdot \dot \varphi) & \text{ak}   \ \dot \varphi \neq 0  \\
        F_{c0} \ \tanh(hc\cdot      \varphi) & \text{ak}   \ \dot \varphi = 0 \quad \text{a} \quad \left|\tau_s\right| \geq F_{c0} \\
        \tau_s &  \text{ak}  \ \dot \varphi = 0 \quad \text{a} \quad \left|\tau_s\right| < F_{c0}
    \end{cases}  \\
    \tau_1 &= F_{c1} \ \dot \varphi
\end{align}
\end{subequations}
a parametre systému sú
\begin{align}
    m_k \qquad I_k \qquad L_T \qquad F_{c0} \qquad hc \qquad F_{c1}
\end{align}
Gravitačná konštanta $g$ je známa.


\section{Zápis pre hľadanie numerického riešenia}

Pre hľadanie numerického riešenia predmetnej pohybovej rovnice je potrebné ju rozložiť na sústavu rovníc prvého rádu. Nech
\begin{subequations}
\begin{align}
    x_1(t) &= \varphi(t) \\
    x_2(t) &= \dot \varphi(t)
\end{align}
\end{subequations}
Potom
\begin{subequations}
\begin{align}
    \dot x_1(t) &= x_2(t) \\
    \dot x_2(t) &= - \frac{1}{I_k}\ m_k \  g \   L_T\  \sin(x_1(t)) - \frac{1}{I_k} \tau_0(t) - \frac{1}{I_k} \tau_1(t)
\end{align}
\end{subequations}
a pritom
\begin{subequations}
\begin{align}
    \tau_s &= - m_k g L_T \sin(x_1) \\
    \tau_0 &= 
    \begin{cases} 
        F_{c0} \ \tanh(hc\cdot x_2) & \text{ak}   \ x_2 \neq 0  \\
        F_{c0} \ \tanh(hc\cdot      x_1) & \text{ak}   \ x_2 = 0 \quad \text{a} \quad \left|\tau_s\right| \geq F_{c0} \\
        \tau_s &  \text{ak}  \ x_2 = 0 \quad \text{a} \quad \left|\tau_s\right| < F_{c0}
    \end{cases}  \\
    \tau_1 &= F_{c1} \ x_2
\end{align}
\end{subequations}

Parameter $hc$ (napr. $hc = 10^6$) a konštantu $g = 9,81$ [m/s\(^2\)] je možné považovať za predvolené. Zvyšné parametre je vo všeobecnosti potrebné určiť.

Hmotnosť ramena kyvadla $m_k$ a vzdialenosť ťažiska od osi otáčania $L_T$ je možné relatívne presne stanoviť meraním. Pre určenie ostatných parametrov ($I_k$, $F_{c0}$, $F_{c1}$) je typicky možné využiť metódy identifikácie systémov na základe nameraných dát z~reálneho systému (časové priebehy výstupu systému).

Uvedený matematicky definovaný model je možné implementovať vo forme funkcie, ktorú vo všeobecnosti využíva ODE solver na hľadanie numerického riešenia. Nasledujúci kód predstavuje implementáciu uvedeného modelu v~jazyku Python (s~využitím knižnice NumPy). 


{\catcode`\-=12
\lstinputlisting[language=Python,
                 caption={Súbor \lstinline{/PY/fcn_dynModel/dynSysFyzKyv_suchvisk.py}},
				 consecutivenumbers=false,
                 ]{../PY/fcn_dynModel/dynSysFyzKyv_suchvisk.py}
}














\section{Verifikácia modelu}

Majme experiment, kde manuálne vychýlime rameno kyvadla laboratórneho zariadenia AeroShield na doraz v smere zápornej hodnoty uhla a následne ho pustíme. Takto získané dáta sú na nasledujúcom obrázku označené ako \emph{meranie}.

% \begin{center}

%     \vbox{%
%         \makebox[\textwidth][c]{%
%         \includegraphics{fj_01_file3_panel_1.pdf}
%         }

%         \makebox[\textwidth][c]{%
%         \includegraphics{fj_01_file3_panel_2.pdf}
%         }

%         \figcaption{ 
%             Voľné kmity ramena kyvadla laboratórneho zariadenia AeroShield. V čase $t = 2,54$ je kyvadlo uvolnené z~vychýlenej polohy.
%         }
%         \label{fj_01_file3}
%     }%vbox

% \end{center}

\begin{center}

    \vbox{%
        \makebox[\textwidth][c]{%
        \includegraphics{fj_02_file3_panel_1.pdf}
        }

        \makebox[\textwidth][c]{%
        \includegraphics{fj_02_file3_panel_2.pdf}
        }

        \figcaption{ 
            Voľné kmity ramena kyvadla laboratórneho zariadenia AeroShield. V čase $t = 2,54$ je kyvadlo uvolnené z~vychýlenej polohy.
        }
        \label{fj_02_file3}
    }%vbox

\end{center}

Od času $t = 2,54$ [s] sú na obrázku zobrazené voľné kmity ramena kyvadla. Verifikácia modelu spočíva v porovnaní nameraných dát a simulácie. Tu sa obmedzíme na grafické porovnanie v rámci obr. \ref{fj_02_file3}, kde sú simulované dáta označené ako \emph{simulácia}.

Pre simuláciu boli použité nasledovné hodnoty parametrov modelu. Podrobnejší opis postupu určenia týchto hodnôt je nad rámec tohto textu.
\begin{align*}
    g &= 9,81 &\ \text{[m/s}^2\text{]} \\
    hc &= 10^6 \\
    m_k &= 12 &\ \text{[g]} \\
    L_T &= 6 &\ \text{[cm]} \\
    I_k &= 0.00007979 &\ \text{[kg m}^2\text{]} \\
    F_{c0} &= 0.00013296 &\ \text{[N m]} \\
    F_{c1} &= 0.00000629 &\ \text{[N m/(rad/s)]}
\end{align*}


\noindent
Model sa považuje za verifikovaný.













% -----------------------------------------------------------------------------

\end{document}

% -----------------------------------------------------------------------------